\subsection{Tether variants}
\label{app:tether}

This appendix defines the nine variants compared in
Fig.~\ref{fig:ablation}(b) and states the calibration protocol that makes them
comparable.

\paragraph{A general form.}
Every variant is an instance of
\begin{equation}
\mathcal{L}_{\mathrm{tether}}(\theta) \;=\; \lambda\,
\mathbb{E}_{s\sim\mathcal{D}}\Big[\,
g\big(\lVert M\,(\pi_\theta(s)-\pi_{\mathrm{ref}}(s))\rVert_2\big)\Big],
\label{eq:tether-general}
\end{equation}
which exposes exactly four free positions: the shape $g$ of the penalty, the
state distribution $\mathcal{D}$ on which proximity is measured, the metric $M$
on action space, and the reference $\pi_{\mathrm{ref}}$ the proposal is pulled
towards. The deployed setting is $g(x)=x^2$, $\mathcal{D}=\mathcal{D}^d_s$,
$M=I$ and $\pi_{\mathrm{ref}}=\pi_s$. Block~(i) of
Fig.~\ref{fig:ablation}(b) varies $g$ and $M$, block~(ii) varies $\mathcal{D}$,
and block~(iii) varies $\pi_{\mathrm{ref}}$. A fifth degree of freedom --
penalising single-step actions versus whole trajectories -- is left untested.

\paragraph{The lift scale $\sigma$.}
The policy head is a deterministic $\tanh$-MLP, so it has no density and no
divergence can be taken between two of its outputs directly. The only way to
speak of a divergence is to lift both policies to isotropic Gaussians
$\mathcal{N}(\pi(s),\sigma^2 I)$ sharing one width $\sigma$. We stress that
$\sigma$ is a property of this lift and not of the policy: nothing in the method
samples from it, no loss estimates it, and it may be chosen freely. Its only
role is to set the scale against which a deviation counts as large, so the
divergences below depend on $\pi_\theta$ and $\pi_s$ only through the
dimensionless ratio $d(s)/\sigma$, with
$d(s) = \lVert\pi_\theta(s)-\pi_s(s)\rVert_2$.

\paragraph{Why the name of the $f$-divergence is not a design choice.}
For two equal-covariance Gaussians every standard divergence is a function of
$d$ alone:
\begin{equation}
\mathrm{KL} = \frac{d^2}{2\sigma^2},\quad
W_2^2 = d^2,\quad
H^2 = 1-e^{-d^2/8\sigma^2},\quad
\mathrm{TV} = 2\Phi\!\Big(\frac{d}{2\sigma}\Big)-1,\quad
\chi^2 = e^{d^2/\sigma^2}-1 ,
\label{eq:divergences}
\end{equation}
with $\Phi$ the standard normal CDF. Two consequences follow, and together they
determine which arms are worth running.

First, all five agree to leading order: each equals $c\,d^2/\sigma^2 +
\mathcal{O}\big((d/\sigma)^4\big)$ for a constant $c$ that $\lambda$ absorbs.
Whenever the proposal stays well inside the lift scale they are the \emph{same}
penalty, and $W_2^2$ is not even approximately but exactly the deployed squared
$L_2$. So ``use KL instead'' or ``use $W_2$ instead'' does not name a different
method here; it names the method we already run, which is why
Fig.~\ref{fig:ablation}(b) labels that row \emph{squared $L_2$ / KL} rather than
listing KL and $W_2$ as separate arms.

Second, where they do differ --- at deviations comparable to or larger than
$\sigma$ --- they differ only through the \emph{shape} $g$ of
Eq.~\ref{eq:tether-general}, and that shape has three qualitative regimes:
sub-quadratic and bounded, so that far-drifted states are progressively released
($H^2$, and $\mathrm{TV}$, which saturates likewise); quadratic, so that every
state keeps pulling in proportion ($\mathrm{KL}$, $W_2$); and super-quadratic,
so that a deviation past some scale is what the loss is about ($\chi^2$, and the
hinge). The experiment therefore varies the shape axis rather than enumerating
names, taking one representative per regime: the saturating penalty for the
first, the deployed squared $L_2$ for the second, and a hard trust region for
the third. We use the hinge rather than $\chi^2$ for the super-quadratic regime
because $e^{d^2/\sigma^2}$ has unbounded gradient and would confound the shape
with an optimisation failure, whereas the hinge realises the same
``free until a threshold, then punished'' behaviour with a bounded curvature and
is the standard trust region of TRPO. $\mathrm{TV}$ is not run separately: it is
a monotone reparametrisation of a saturating penalty and would duplicate $H^2$.

What is \emph{not} interchangeable is the second-order structure. For a Gaussian
policy with fixed $\sigma$ the Fisher information is $F = G/\sigma^2$, where $G$
is the Gauss--Newton matrix of the policy output on $\mathcal{D}^d_s$, so the
deployed penalty is a KL trust region measured in the correct Fisher metric.
That is the metric $M$ of Eq.~\ref{eq:tether-general}, and it is what the
\emph{parameter-space $L_2$} and \emph{weight decay} rows give up.

\paragraph{(i) The shape $g$ and the metric $M$.}
Let $d(s) = \lVert \pi_\theta(s)-\pi_s(s)\rVert_2$.
\begin{itemize}\itemsep2pt
\item \textbf{Squared $L_2$ / KL (ours).} $g(x)=x^2$, i.e.
      $\lambda\,\mathbb{E}_s[d(s)^2]$. Quadratic in the deviation, so every
      state contributes and large deviations dominate.
\item \textbf{Saturating (Hellinger$^2$).}
      $g(x) = 1-e^{-x^2/\tau}$, i.e. $\lambda\,\mathbb{E}_s[1-e^{-d(s)^2/\tau}]$,
      exactly $H^2$ from Eq.~\ref{eq:divergences} with $\tau = 8\sigma^2$; the
      calibrated $\tau$ is thus a statement about where the saturation knee
      sits, and $\tau\to\infty$ recovers the deployed squared penalty. States
      that have already drifted far are progressively released, which makes
      this arm the mirror image of the hinge on the $g$ axis.
\item \textbf{Trust region (hinge).} $g(x) = \mathrm{relu}(x-\epsilon)^2$, i.e.
      $\lambda\,\mathbb{E}_s[\mathrm{relu}(d(s)-\epsilon)^2]$: a soft-penalty
      realisation of the hard constraint $\mathrm{KL}\le\delta$ of TRPO, with
      $\epsilon = \sigma\sqrt{2\delta}$. Deviations below $\epsilon$ are free.
      The two hinge rows share this form exactly and differ only in the radius
      of that free zone, but the radius decides what the penalty is for.
      With $\epsilon=0.05$, comparable to the working point itself, most states
      still sit outside the free zone and the arm behaves like the deployed
      quadratic penalty, which is what it scores. With $\epsilon=0.20$, several
      times the working point, almost every state is free and the penalty has
      no grip on the bulk of the distribution; calibration compensates by
      raising $\lambda$ from $0.3$ to $100$, so the loss becomes zero for
      typical states and enormous for the few beyond the radius. It stops
      tethering a distribution and starts clipping outliers, and the diagnostics
      say so: it is the only variant that never quite reaches the working point
      ($D_{\mathrm{src}}=0.031$ against $0.025$--$0.027$ elsewhere), it pays
      $1.9\times$ the cloning loss of the tight setting ($0.026$ vs $0.014$),
      and it produces the \emph{thinnest} deviation tail of any arm --- which,
      together with its $11$-point deficit, is one of the observations that
      rules out tail shape as the causal channel.
\item \textbf{Parameter-space $L_2$.} $c\,\lVert\theta-\theta_s\rVert_2^2$,
      where $\theta_s$ are the stage-1 policy-head weights. This is the same KL
      trust region with the Fisher metric replaced by the identity. It is the
      only row that changes $M$ rather than $g$, and it is expensive: the
      fine-tuning update lies almost entirely in directions to which the policy
      output is insensitive, so an isotropic parameter penalty cannot tell those
      directions apart and has to shrink all of them, paying roughly
      $3.4\times$ the cloning loss for the same proximity.
\item \textbf{Weight decay.} Adam \texttt{weight\_decay} on the policy head.
      This is not a divergence at all: its attractor is $\theta=0$, not
      $\theta_s$. As the coefficient grows the policy is driven towards
      $\pi_\theta\equiv 0$, whose distance to the source policy is
      $\mathbb{E}_s\lvert\pi_s(s)\rvert$, and the achievable proximity bottoms
      out an order of magnitude short of the working point before turning back
      up. Generic capacity control cannot produce a policy close to $\pi_s$;
      naming $\pi_s$ as the reference is necessary.
\end{itemize}

\paragraph{(ii) The state distribution $\mathcal{D}$.}
The deployed setting measures proximity on all of $\mathcal{D}^d_s$, the same
source data the world model was fit on. The variant restricts it to the steps
that carry both an expert flag and a success flag ($14.6\%$ of transitions) --
the subset on which $\pi_s$ itself was cloned, and therefore the subset where
the recorded actions are closest to $\pi_s$.

\paragraph{(iii) The reference $\pi_{\mathrm{ref}}$.}
The deployed setting pulls towards the frozen source policy $\pi_s$ evaluated at
$s$. The variants replace it with the action recorded in the source dataset at
that transition, either over all of $\mathcal{D}^d_s$ or over the
expert\,$\wedge$\,success subset. Recorded actions carry a systematic offset
from $\pi_s$ -- an elementwise $0.170$ over the full dataset -- which no
coefficient removes, so that arm cannot reach the working point at any
$\lambda$. Filtering to expert\,$\wedge$\,success cuts the offset to $0.025$ and
does restore the working point, which is why the two rows in block~(iii) fail
for different reasons.

\paragraph{Calibration.}
A larger $\lambda$ trivially buys proximity to $\pi_s$, so comparing the
variants at a common coefficient would compare strengths rather than forms.
Before any target interaction we therefore run stage~2 alone for each variant
and sweep its knob until it matches the deployed arm on
\begin{equation}
D_{\mathrm{src}} \;=\; \frac{1}{d_a}\,
\mathbb{E}_{s\sim\mathcal{D}^d_s}\big\lVert
\pi_\theta(s)-\pi_s(s)\big\rVert_1 ,
\label{eq:dsrc}
\end{equation}
the elementwise mean absolute deviation from the source policy over source
states, with $d_a$ the action dimension. Stage~3 is only run once a variant sits
at $D_{\mathrm{src}}\approx0.025$. This is what makes the resulting
$\lambda$ values span four orders of magnitude in
Fig.~\ref{fig:ablation}(b): they are the output of the calibration, not a
choice. Two variants -- weight decay and the unfiltered recorded actions --
cannot reach the working point at any strength and are reported at their
attainable floor; for those rows the floor \emph{is} the result.
